(1) 忽略位移电流时 \(\nabla\times\mathbf H=\sigma\mathbf E\)、\(\nabla\times\mathbf E=-\partial_t\mathbf B=-\mu\partial_t\mathbf H\)。取旋度并用 \(\nabla\cdot\mathbf H=0\) 得 \(\nabla^2\mathbf H=\mu\sigma\,\partial_t\mathbf H\)。对 \(e^{-i\omega t}\) 谐波，\(H_z''=-i\mu\sigma\omega H_z\)，衰减根为 \(q=(1-i)/\delta\)，其中 \(\delta=\sqrt{2/(\mu\sigma\omega)}\)。所以 \(\boxed{H_z(x)=H_0e^{-(1-i)x/\delta}}\)，幅度按 \(e^{-x/\delta}\) 衰减，场相位沿深度正向滞后于表面；不能选增长的根。

(2) 由于 \((\nabla\times\mathbf H)_y=-H_z'\)，有 \(j_y=-H_z'=qH_z\)，\(E_y=j_y/\sigma=qH_z/\sigma\)。因此 \(\boxed{Z_s=(1-i)/(\sigma\delta)}\)。在本题 \(e^{-i\omega t}\) 约定下 \(\operatorname{Im}Z_s<0\)；若改用 \(e^{+i\omega t}\)，电抗虚部应相应反号，而可测耗散实部 \(\operatorname{Re}Z_s=1/(\sigma\delta)\) 不变。表面阻抗单位为欧姆，不能与体电阻率 \(1/\sigma\) 混淆。

(3) 切向 \(\mathbf E=E_y\hat y,\mathbf H=H_z\hat z\)，所以 \(\mathbf E\times\mathbf H^*\) 指向导体内部 \(+\hat x\)。平均入射功率密度为 \(P_A=\tfrac12\operatorname{Re}(E_yH_z^*)|_{0}=\frac{|H_0|^2}{2\sigma\delta}\)。体焦耳热积分为 \(\int_0^\infty(\sigma/2)|E_y|^2dx=\frac{\sigma}{2}\frac{2|H_0|^2}{\sigma^2\delta^2}\int_0^\infty e^{-2x/\delta}dx=\frac{|H_0|^2}{2\sigma\delta}\)。两法完全一致，表明稳态周期内储存磁能的净变化为零。

(4) 从 \(\delta\propto\omega^{-1/2}\) 得 \(\operatorname{Re}Z_s\propto\sqrt\omega\)，固定表面磁场下平均损耗 \(P_A\propto\sqrt\omega\)。形式 \(\omega\to0^+\) 时 \(\delta\to\infty,Z_s\to0\)，但无限半空间的有限频边界问题与真正静态磁场在整个导体中穿透的极限不可直接交换。频率升高到 \(\omega\varepsilon\) 与 \(\sigma\) 可比时位移电流不可忽略；当 \(\delta\) 与电子平均自由程或表面粗糙尺度相当时局域Ohm定律及平滑边界假设也会失效。
